\section{Stakeholder Analysis}
The needs of stakeholders are analyzed by comparing the as-is state of their capabilities against a to-be (desired) state derived from a review of external references and the state-of-the-art.
This is a preliminary discussion that will be later formalized into verifiable requirements.
The analysis characterizes the needs, requirements, interfaces and technical criteria of nanosatellite EPS, both as observed within ASCEND and in the consulted references, and thus represents the deliverable of the first objective of this project.

\subsection{As-is capabilities}
For the CUBICS mission, a collaboration between ASCEND researchers and members of a student club at ÉTS is being proposed.
The present work is aimed at serving them both.

ASCEND researchers are familiar with systems engineering and even MBSE, but the nature of their position does not allow them to contribute significant time to the practical aspect of CubeSat design and development.
In contrast, ASCEND-associated students who wish to participate in CUBICS have years of experience in building systems for student engineering competitions but follow no systems engineering method.

The students have built CubeSat-format payloads for rocketry competitions, for which they verified payload power requirements only after fabrication.
This verification was observed to take the form of reusing a legacy battery pack for which the original calculations have been lost, connecting all loads, and performing mission operations repeatedly until the stored energy was depleted.
If the time until energy depletion was greater than the duration of one operations cycle as defined by the competition rules, the power requirements were taken as verified.
These experiments were reported to result in battery lifespans several times greater that one operations cycle, with the students opining that the battery pack was oversized.

In contrast, the student club members have ample experience in verifying requirements related to rocket mass, maximum height, and maximum speed before fabrication, via inspection and analysis in simulation.
These requirements are given as numerical targets by each rocketry competition, or derivated from the top-place results of previous competitions.
These competition requirements are always expressed in terms of measurable quantities, which is essential for competition grading.

ASCEND researchers report an EPS design that has been in progress for nine months.
Unclear definition of requirements has been a recurring issue, and it has forced redesigns when it is discovered that the existing requirement list is incomplete.
Requirement changes force the researchers to wait at least one week, until the next one hour-long weekly meeting, to assess their impact, because the traceability from requirements to design has to be established manually each time.

\subsection{Stakeholder concerns}
\label{sec:stakeholder-concerns}

Following OOSEM, stakeholder discussions identified concerns about the cost and effort of closing capability gaps.
These concerns inform framework requirements and measures of performance.

A project manager and systems engineer at ASCEND is concerned with maintaining verifiable requirements, redesign after fabrication, and tools suited to the team's mission scale and resources \citep[Secs.~4--5]{handpicked-0003}\citep[Sec.~1]{review-0573}\citep[Sec.~5.1]{buruso2026modelbased}.
That stakeholder is also concerned with requirement traceability, communicating verification and validation evidence, and updating EPS and other subsystem power requirements as designs change.
Maintaining consistency across requirements, models and verification evidence also demands effort \citep[Secs.~4.4 and~5.1]{buruso2026modelbased}\citep[Secs.~3--4]{review-0853}\citep[Sec.~V]{citation-0136}.

An EPS engineer is concerned with repeating verification after requirement changes, tracing design decisions, revising designs after failed verification, updating component requirements and communicating designs during technical reviews.
Further concerns include verifying safe and effective charging and power management \citep[pp.~1--3 and~6]{handpicked-0001}\citep[Sec.~3]{handpicked-0003}\citep[Secs.~5.1 and~6]{review-1251}.
Engineers responsible for other subsystems share the concern about verification cost and effort.

Taken together, the aforementioned as-is capability gaps and associated concerns can be formulated as a problem that motivates the current project:

\begin{quote}
Evidence fragmentation during the early lifecycle phases of nanosatellite electrical power systems hinders the traceability between needs and requirements, architecture decisions, and the analytical evidence that supports them. Furthermore, it hinfers the reproducibility of analyses and the capability to detect design inconsistencies before advacing to physical verification and validation phases.
\end{quote}

\subsection{Potential to-be capabilities}
After observing as-is stakeholder capabilities and recording them objectively, the OOSEM method suggests analyzing the capabilities of existing external enterprises and systems. From this, capability gaps or the lag behind the state-of-the-art can be identified and later formalized into the requirements for the systems and processes that will improve these capabilities, in this case the proposed framework.
Sources for this comparison include the academic literature, crucial for ASCEND as a research group, and normative documents.
Said sources were discussed amply in Chapter \ref{ch:marco}, to which the reader is referred in order to avoid repetition.
